\section{Introduction}
\label{sec:intro}

A wave of lunar exploration is placing new demands on the geodetic
infrastructure of the Moon. ESA's Moonlight programme and its Lunar
Communications and Navigation Service (LCNS), NASA's LunaNet and Lunar
Communications Relay and Navigation Systems, and Japan's Lunar Navigation
Satellite System will together provide positioning, navigation, and timing (PNT)
for surface and orbital users \citep{israel2020lunanet,fienga2024moonlight}. Every
one of these services is expressed in, and only as good as, a common lunar
reference frame. Recognizing this, the community has begun to define and realize
the International Lunar Reference Frame (ILRF): a principal-axis system co-rotating
with the Moon, with origin at the lunar centre of mass \citep{sosnica2025ilrf}.

The first ILRF realization exposes a sharp limitation. \citet{sosnica2025ilrf}
report that a poor spatial distribution of the lunar laser-ranging (LLR)
retroreflector network---five near-side reflectors, clustered near the equator,
with a single station in the southern hemisphere---induces a correlation as large
as $\corr \approx -0.97$ between the lunocentre $X$-translation and the overall
frame scale. This near-degeneracy dominates the origin error budget: the combined
solution carries a mean position error of $17.6$~cm over 2010--2030, of which
$15.3$~cm comes from the origin. In terrestrial terms this is the lunar analogue
of the origin-and-scale problem that the terrestrial reference frame community has
spent decades taming through geometric and technique diversity
\citep{altamimi2016itrf2014}: the datum---the seven degrees of freedom (three
translations, one scale, three rotations) that no internal measurement can
fix---must be pinned down by observations whose geometry actually senses those
degrees of freedom.

The empirical $-0.97$ is a symptom. The design questions behind it are not yet
answered: \emph{Why} is it the $X$-translation and the scale, specifically, that
lock together? \emph{Which} additional observations would separate them, and
which would not? And, under a finite budget, \emph{what is the cost-optimal way}
to buy the geometry that breaks the degeneracy? These are identifiability and
experiment-design questions, and they are the subject of this paper.

\paragraph{Related work.}
Existing lunar-PNT performance analyses are, almost without exception, analyses of
the \emph{user} positioning problem: given a realized frame and a set of beacons,
how accurately can a receiver be located. \citet{pohlmann2025hybrid} derive
realistic error models and Cram\'er--Rao lower bounds for hybrid satellite and
cooperative navigation, showing sub-metre user accuracy with as few as two visible
satellites when aided by a reference station; open simulators such as LuPNT
\citep{iiyama2023lupnt} target the same user-side problem. These works take the frame as given. The
frame-\emph{realization} problem---the identifiability of the datum parameters
themselves---has been addressed empirically, as a by-product of a specific
inversion \citep{sosnica2025ilrf}, and reviewed qualitatively:
\citet{fienga2024moonlight} note that additional retroreflectors and orbiter
altimetry could improve frame accuracy, but do not quantify which geometry
resolves which degeneracy. To our knowledge, no prior work provides a general
identifiability classification of the internal-ranging lunar datum, nor a
cost-aware design of the observations that break its dominant degeneracy. This
paper fills that gap. It is the geodetic-datum dual of the user-side Cram\'er--Rao
analyses: the same Fisher-information machinery, applied to the seven datum
parameters rather than to a user state.

\paragraph{Contributions.}
We make three contributions, each accompanied by an explicit statement of what is
proven versus what is representative.

\begin{enumerate}
  \item \textbf{A null-space classification theorem} (Section~\ref{sec:classification})
  for the internal-ranging lunar datum. We show that (i) each scalar range
  observation contributes rank one to the seven-parameter datum Fisher
  information, so a single observation leaves a six-dimensional null space;
  (ii) for a near-side cluster under a fixed body orientation the origin-$X$ and
  scale columns are near-collinear, placing their difference in (or arbitrarily
  near) the null space; and (iii) physical libration lifts the datum defect to
  zero while leaving the pair near-degenerate, with the marginal information a
  monotone function of librational excursion and arc length. The classification
  is a geometric property and is engine-verified.

  \item \textbf{A radial-ambiguity characterization} (Section~\ref{sec:multitechnique}).
  Building a unified datum-partials substrate for LLR, two-station lunar VLBI
  differential delay, and orbiter-to-beacon range over the full seven-parameter
  Helmert datum, we show that the lunocentre-$X$/scale degeneracy is a
  \emph{radial} ambiguity. Radial, depth-diverse ranging (far-side orbiter
  tracking) breaks it directly; transverse VLBI helps only indirectly, because its
  differential-delay information lands in the transverse translation, lifting the
  degenerate pair only through the operator-monotonicity of the marginal
  information.

  \item \textbf{A cost-aware design law} (Section~\ref{sec:oed}). We cast frame
  realization as gauge-constrained optimal experiment design, verify a greedy
  design heuristic against an exact brute-force oracle, and compute a
  cost/degeneracy Pareto frontier over a representative multi-technique menu. A
  far-side orbiter range delivers $\approx\!8\times$ the raw degeneracy-breaking
  gain of a limb VLBI baseline for comparable exposure ($\approx\!4.8\times$ per
  unit cost), making depth diversity the priority acquisition in any
  budget-constrained realization campaign.
\end{enumerate}

\paragraph{Honesty scope.}
Throughout, we separate two kinds of claim. The \emph{structure}---the rank
count, the sign and near-collinearity of the degenerate directions, the
classification, the improvement \emph{direction}, and the correctness of the
decomposition linear algebra---is proven or independently validated. The
\emph{magnitudes}---the residual correlation, the Cram\'er--Rao bounds, the
degeneracy-metric values, and the relative costs---are computed from
\emph{representative} lunar geometry, noise, and economics; they are labelled
\Modelled{} (and tagged \modtag) and are \emph{not} claims about any real mission
or a reproduction of the $-0.97$/$15$~cm magnitudes of \citet{sosnica2025ilrf}.
Section~\ref{sec:validation} makes this firewall precise. Every number in the
paper is regenerated from the committed engine by a single command
(Section~\ref{sec:availability}).
